% theme: moon_and_planets
\MajorQuestion{月と惑星}
\SetRowSpaceQanda

\Instruction{次の問いに答えなさい。}
\QQ{月が地球のまわりを回ることを\NamedBlank{moon_revolution}という。}{公転}
\QQ{太陽が月によってかくされる現象を\NamedBlank{solar_eclipse}という。}{日食}
\QQ{月が見えない形を\NamedBlank{new_moon}という。}{新月}
\QQ{月全体が丸く明るく見える形を\NamedBlank{full_moon}という。}{満月}
\QQ{月の見かけの形が変化することを\NamedBlank{moon_phase}という。}{満ち欠け}

\Instruction{\strut}
\QQ{惑星のまわりを回る天体を\NamedBlank{satellite}という。}{衛星}
\QQ{月が地球のまわりを1回\RefBlank{moon_revolution}する期間はどのくらいか。}{約27日}
\QQ{\RefBlank{solar_eclipse}のうち、太陽が完全にかくされるものを何というか。}{皆既日食}
\QQ{\RefBlank{new_moon}から約1週間後に見える半月を何というか。}{上弦の月}
\QQ{\RefBlank{full_moon}から約1週間後に見える半月を何というか。}{下弦の月}

\Instruction{\strut}
\QQ{月が1日に1回、東から西へ動くように見える運動を\NamedBlank{moon_diurnal}という。}{日周運動}
\QQ{月の表面に見える、丸いくぼみを何というか。}{クレーター}
\QQ{月が\RefBlank{moon_revolution}と同じ周期で自ら回転することを何というか。}{自転}
\QQ{月が地球の影に入る現象を何というか。}{月食}
\QQ{\RefBlank{new_moon}から次の\RefBlank{new_moon}までの期間は約何日か。}{約30日}

\Instruction{\strut}
\QQ{夕方、西の空に見える金星を何というか。}{よいの明星}
\QQ{同じ時刻に見える月の位置は、毎日少しずつどちらからどちらへ移るか。}{西から東}
\QQ{月の\NamedBlank{moon_diurnal}では、地球から見て月はどの方角からどの方角に動いていくように見えるか。}{東から西}
\QQ{月の光は、何の光が反射されたものか。}{太陽}
\QQ{明け方、東の空に見える金星を何というか。}{明けの明星}
